\documentclass{article}

\usepackage[utf8]{inputenc}     % pdfLaTeX
\usepackage[T1]{fontenc}

\usepackage{amsmath}
\usepackage{amsthm}
\usepackage{amssymb}

\title{$\sqrt{2}$}
\author{Rovenszky Robin}
\date{Last updated: 13 March 2026}

\newtheorem{theorem}{Theorem}
\newtheorem*{remark}{Remark}

\begin{document}

\maketitle

\section{What is $\sqrt{2}$?}
We understand that it must exist, since in a right trinagle with legs $a = 1, b= 1, \alpha = 90$, where $\alpha$ is their enclosed angle, the hypothenuse has length $\sqrt{2}$. So what is $\sqrt{2}$?
\\
\begin{theorem}
    For an x satisfying $x^2 = 2$, x may also satisfy $x = \frac{p}{q}$ with $p, q \in \mathbb{Z}$, $gcd(p, q) = 1$
\end{theorem}
\begin{proof}
    Assume the statement is true. Then
    \begin{equation}
        \frac{p^2}{q^2} = 2 \implies p^2 = 2q^2 \implies 2 \mid p^2 \implies 2 \mid p \implies 4 \mid p^2 \implies
    \end{equation}
    \begin{equation}
        \implies 4 \mid 2q^2 \implies 2 \mid q^2 \implies 2 \mid q \implies \gcd(p, q) \ge 2 \Rightarrow\!\Leftarrow
    \end{equation}
    Or alternatively, we could have also finished the proof via
    $p^2 = 2q^2$. Now we notice that $p^2$ has two in its prime factorization to an even power, while $2q^2$ has two in its prime factorization to an odd power (since $\gcd(p,q) = 1$).
\end{proof}

We can generalize this result to show that for all $a \in \mathbb{Z}_{\ge 0}$, if $x^2 = a$, then if $a$ contains a prime factor to an odd exponent (so $a$ is not a perfect square), then and only then $x \in \mathbb{R} \setminus \mathbb{Q}$.

\section{Finding $\sqrt{2}$}
Since we cannot find irrational numbers, we will approximate it. Let us start with:
\[
1 < x < 2 \Leftrightarrow 1 < x^2 < 4
\]
\begin{theorem}
    \[
    0 < x < y \implies 0 < x^2 < y^2
    \]
\end{theorem}
\hrule
\vspace{1em}

Before we prove our theorem, let us write out the axioms of order (also known as axioms of arrangement):
\begin{enumerate}
    \item Axiom of Trichotomy: Exactly one of the following can be true: $a < b, b < a, a = b$.
    \item Addition preserves order: $a < b$ and $c$ arbitrary $\implies$ $a + c < b + c$.
    \item  Axiom of Transitivity: $a < b$ and $b < c$ $\implies$ $a < c$.
    \item Multiplication by positive preserves order: $0 < a$ and $b < c$, then $ab < ac$.
    \item Positivity of 1: $0 < 1$.
\end{enumerate}
We will use the distributivity of addition (or that addition preserves order) to prove the following:
\[
    x < 0 \implies 0 < -x
\]
\begin{proof}
\[
    x + (-x) < 0 + (-x)
\]
\[
    0 < -x
\]
\end{proof}
Now we use the distributivity of multiplication (or that multiplication by positive preserves order) to prove the following:
\[
    (-1)\cdot x = -x
\]
\begin{proof}
\[
    (1 + (-1)) \cdot x = 0
\]
\[
    1 \cdot x + (-1) \cdot x = 0 \implies (-1) \cdot x = -x
\]
\end{proof}
For this however, we need the following statement aswell, for which we will again use the distributivity of multiplication:
\[
    0\cdot x = 0
\]
\begin{proof}
\[
    (0+0) \cdot x = 0 \cdot x
\]
\[
    0 \cdot x + 0 \cdot x = 0 \cdot x
\]
\[
    0 \cdot x = 0
\]
\end{proof}

\hrule
\vspace{1em}

Now we proceed with a proof of Theorem 3 using the distributivity of multiplication:

\begin{proof}
    \begin{equation}
    0 < x < y \implies 0 < x^2 < y^2
    \end{equation}
    Then multiplying by x yields:
    \[
    x^2 < xy
    \]
    and multiplying (27) by y gives us:
    \[
    xy < y^2
    \]
    Combining the latter two yields:
    \[
    x^2 < xy < y^2 \implies x^2 < y^2
    \]
\end{proof}

\begin{remark}
    We can also prove that $y < x < 0 \implies 0 < -x < -y$ using the distributivity of multiplication.
\end{remark}

\begin{proof}
    \[
    (-x) \cdot (-x) < (-y) \cdot (-x)
    \]
    \[
    x^2 < xy < y^2
    \]
\end{proof}

\end{document}